\begin{tikzpicture}[
    scale=0.6,
    transform shape,
    font=\small,
    block/.style={
        draw,
        rounded corners=2pt,
        thick,
        minimum width=3.1cm,
        minimum height=1.05cm,
        align=center
    },
    src/.style={block, fill=gray!12},
    trans/.style={block, fill=blue!20},
    met/.style={block, fill=red!12},
    agg/.style={block, fill=orange!25},
    sink/.style={block, fill=green!15},
    arrow/.style={-{Stealth[]}, thick},
    art/.style={font=\scriptsize\itshape, text=gray!60!black, align=center}
]

% ============================================================
% GLOWNY PRZEPLYW
% ============================================================
\node[src]   (zrodlo) at ( 0.0, 0) {Źródła danych\\[-2pt]\scriptsize{typowane pola}};
\node[trans] (przeksz) at ( 6.2, 0) {Węzły\\[-2pt]przekształceń};
\node[met]   (metryki) at (12.4, 0) {Węzły metryk};
\node[agg]   (agreg)   at (18.2, 0) {Agregacja};
\node[sink]  (raport)  at (23.6, 0) {Raport};

\draw[arrow] (zrodlo)  -- node[art, above] {artefakty} (przeksz);
\draw[arrow] (przeksz) -- node[art, above] {artefakty} (metryki);
\draw[arrow] (metryki) -- node[art, above] {wyniki\\elementowe} (agreg);
\draw[arrow] (agreg)   -- (raport);

% ============================================================
% CO KRYJE SIE POD WEZLAMI PRZEKSZTALCEN
% ============================================================
\node[align=center, font=\scriptsize, text=gray!60!black, below=1.05cm of przeksz]
    (opis) {transkrypcja, synteza mowy, embedding,\\
            indeksowanie, wyszukiwanie, zawężanie,\\
            augmentacja, generacja odpowiedzi};
\draw[gray, dashed] (przeksz.south) -- (opis.north);

\node[align=center, font=\scriptsize, text=gray!60!black, below=1.05cm of metryki]
    (opism) {dostawiane automatycznie wszędzie,\\gdzie ich wejścia są obecne};
\draw[gray, dashed] (metryki.south) -- (opism.north);

% ============================================================
% TABLICA ARTEFAKTOW JAKO WSPOLNE PODLOZE
% ============================================================
\node[draw=gray, dashed, rounded corners=4pt, inner sep=0.42cm,
      fit=(zrodlo)(przeksz)(metryki)(agreg)(raport),
      label={[font=\scriptsize\itshape, gray]above:wspólna tablica artefaktów przebiegu}] {};

\end{tikzpicture}
